\documentclass{article}
\usepackage[margin=2cm]{geometry}
\usepackage{amsmath}
\usepackage{graphicx}
\usepackage{booktabs}
\usepackage[utf8]{inputenc}
\usepackage{ragged2e}
\setlength{\emergencystretch}{3em}
\begin{document}
\sloppy

\noindent\hfill\includegraphics[height=1.8em]{assets/figure_header_logo.png}\par\noindent\hrule\vspace{4pt}

We can rewrite W with projector operators


HS theory can also be considered in the sense of
AdS/CFT correspondence, where the HS theory of massless HS fields corresponds to a limiting case of the string.
theory for the string tension going to zero. Large N superconformal field theories were studied as holographic.
duals for higher spin gauge theories in perturbative expansion around AdS spacetime . Klebanov and Polyakov
have proposed a duality between the singlet sector of.
the critical 3-d O(N) vector model with $\{ \phi ^ { \mathbf { a } } \phi ^ { \mathbf { a } } \} ^ { \mathbf { z } }$ interaction and minimal bosonic theory in AdS4 which contains
massless gauge fields with even spin. The analog of AdS4
conjecture by Klebanov and Polyakov appeared in AdS3
conjecturing a duality between a complex scalar coupled
to higher-spin fields in Vasiliev's gravity in 3 dimensions
and WN minimal model CFT in t'Hooft limit denoted
by coset representation


where A are Chern-Simons gauge fields which take value
in the Lie algebra$\mathbf { h s } [ \lambda ]$ In this formulation the equation


$\psi$ elements of the Clifford pairs such that $W ^ { * } ~ =$
$- \mathcal { P } . { + } A - \mathcal { P } . { - } A$for


gives


\begin{equation}
\mathcal { P } _ { \pm } \psi _ { 1 } = \psi _ { 1 } \mathcal { P } _ { \pm } = \pm \mathcal { P } _ { \pm } ~
\end{equation}


\begin{equation}
d W = W \wedge \star W \qquad \qquad ( 6 )
\end{equation}


\begin{equation}
\mathcal { P } _ { \pm } = \frac { 1 \pm \psi } { 2 } \qquad ( 4 )
\end{equation}


\begin{equation}
\begin{array} { r l r } & { } & { \quad \mathrm { s c o n t e m a } } \\ & { } & { \quad ~ }\end{array}
\end{equation}


\begin{equation}
\begin{array} { c } { d A + A \wedge \star A = 0 \qquad \qquad ( 7 ) } \\ { D \overline { \mathbf { A } } + \overline { \mathbf { A } } \wedge \star \overline { \mathbf { A } } = 0 \qquad \qquad \qquad ( 8 ) } \end{array}
\end{equation}


where we define the t'Hooft limit with $N, k ~ \rightarrow ~ \infin$ for
$\begin{array} { r } { \lambda \equiv \frac { N } { k + N ^ { * } } } \end{array}$ This duality has been verified by the number
o1 stuuies, correspondence of global symmetries in bulk
and at the boundary , correspondence of the 1-loop partition function in the bulk and at the large NCFT , and
partition function of the HS black hole at high temperature in the bulk and at the boundary CFT, as well as
for the 3-point functions when $\begin{array} { r } { \lambda = \frac 1 2 } \end{array}$ and ${ \mathfrak { s } } = 2, 3, 4$ for
the scalar-scalar-HS field (00s) corretator in the tHooft
limit. Via three-point functions, tests of the conjecture
have been done in (for O0s correlator with general $\lambda$
and in . In this work, we extract the coupling of the
00s three-point correlator by considering the linearised
Vasiliev equations of motion, and we verify it by choosing the spin to be three and comparing with result in
. The result corresponds to coupling of the three-point
correlation function up to selected normalisation. While
we consider general$\lambda$ the similar work has been done
for the fixed in . The structure of the work is as fol-$\lambda$
lows: In the section two we consider linearised equations
of motion in the Vasiliev's theory, in the section three we
consider the higher spin field in the metric formulation,
and in section four we conclude.


are specific coefficients dependent $\lambda$ and defined ac
cording to conventions The equations describe interaction of arbitrary higher spin background with lienarized
scalars. The coupling that we are interested in can be extracted from rewriting the master field B as a linearized
fluctuation around vacuum value $\nu$


where A and $A$ are positive polynomials of the positive
degree in products of deformed oscillators. () and ( are
in that case equal to field equations $h s [ \lambda ] \otimes h s [ \lambda ]$ Chern-
Simons theory.
The generators of.$\ln s [ \lambda ]$are defined with spin index$\pmb { s }$ and
mode index $m$ as $V . { i n } ^ { \ast }$ for$s \geq 2$ while $| m | < s$ and obey
star product


where


\begin{equation}
V _ { m } ^ { s } \star V _ { n } ^ { t } = \sum _ { u = 1, 2, 3 } ^ { s + t - | s - t | - 1 } g _ { u } ^ { s t } ( m, n ; \lambda ) V _ { m - n } ^ { s + t - n } \qquad ( 9 )
\end{equation}


\begin{equation}
g _ { u } ^ { s t } ( m, n ; \lambda ) = ( - 1 ) ^ { u + 1 } g _ { n } ^ { t s } ( m, n ; \lambda ) ~
\end{equation}


\subsection*{I. LINEARIZED EQUATIONS OF MOTION}


\begin{equation}
B = \nu + \mathcal { P } _ { + } \psi _ { 2 } C ( x, \bar { y } _ { \alpha } ) + \mathcal { P } _ { - } \psi _ { 2 } \hat { C } ( x, \bar { y } _ { \alpha } ) ~ ~ ~ ~ ~ ~ ( 1 1 )
\end{equation}


Let us first consider the coefficient coming from the
Vasiliev linearised equations. Vasiliev's theory contains
five equations for the master fields W which is spacetime
1-form, B and $s . { r }$which are spacetime O-forms. The generating functions are dependent on the coordinates of the
spacetime, auxiliary bosonic twistor variables (referred
to as "oscillators) and Clifford element pairs, where in
definitions we follow conventions from
and various other ingredients are used to define the "deformed oscillator star-commutation relations which give
rise to$h \mathbf { s } [ \lambda ]$ higher spin algebra. Two of the above mentioned equations that will be of the interest here are


and expanding the master field$C$in the deformed oscil-.
lators$\bar { y } . { \alpha }$in the equation


That allows us determining the generalised Klein-Gordon
(KG) equation in the background of HS fields. While the
expansion of the master field C in formalism of bosonic
Vasiliev's theory is given by


\begin{equation}
C = C _ { 0 } ^ { 1 } + C ^ { \alpha \delta } \tilde { y } _ { \alpha } \tilde { y } _ { j i } + C ^ { \alpha \beta \sigma \lambda } \tilde { y } _ { i \cdot } \tilde { y } _ {, i } \tilde { y } _ { \sigma } \tilde { y } _ { \lambda } + \dots \qquad ( \mathbf { 1 3 } )
\end{equation}


\begin{equation}
d C + A \star C - C \star \overline { A } = 0. \qquad \qquad ( 1 2 )
\end{equation}


\begin{equation}
\begin{array} { c } { \cdots \cdots \qquad \cdots \cdots \cdots \cdots \cdots \cdots \cdots \cdots \cdots \cdots } \\ { d W = W \wedge \star W \qquad \qquad \qquad \qquad ( 2 ) } \\ { d B = W \star B - B \star W \qquad \qquad \qquad \qquad ( 3 ) } \end{array}
\end{equation}

\end{document}
